\documentclass[11pt,a4paper]{article}

% ---- Packages (all standard on Overleaf; compiles with pdfLaTeX) ----
\usepackage[margin=2.3cm]{geometry}
\usepackage{xcolor}
\usepackage{enumitem}
\usepackage{booktabs}
\usepackage{listings}
\usepackage{tikz}
\usetikzlibrary{arrows.meta, positioning, shapes.geometric, fit, backgrounds}
\usepackage[colorlinks=true, linkcolor=blue, urlcolor=blue]{hyperref}
\usepackage{titlesec}

% ---- Colors ----
\definecolor{ink}{HTML}{1F2937}
\definecolor{accent}{HTML}{2563EB}
\definecolor{agent}{HTML}{7C3AED}
\definecolor{soft}{HTML}{EEF2FF}
\definecolor{code}{HTML}{F3F4F6}
\definecolor{green}{HTML}{059669}

% ---- Code listing style ----
\lstdefinestyle{plain}{
  backgroundcolor=\color{code},
  basicstyle=\ttfamily\small,
  breaklines=true,
  frame=single,
  framesep=6pt,
  rulecolor=\color{code},
  xleftmargin=6pt,
  showstringspaces=false
}
\lstset{style=plain}

\titleformat{\section}{\Large\bfseries\color{ink}}{\thesection.}{0.5em}{}
\titleformat{\subsection}{\large\bfseries\color{accent}}{\thesubsection}{0.5em}{}

% ---- TikZ node styles ----
\tikzstyle{step}=[rectangle, rounded corners=3pt, draw=accent, fill=soft,
  text width=8.4cm, align=left, inner sep=7pt, font=\small]
\tikzstyle{ai}=[rectangle, rounded corners=3pt, draw=agent, fill=agent!8,
  text width=8.4cm, align=left, inner sep=7pt, font=\small]
\tikzstyle{arrow}=[-{Stealth[length=3mm]}, thick, draw=ink]

\begin{document}

% ======================= TITLE =======================
\begin{center}
  {\Huge\bfseries\color{ink} The Dependency EOL Firewall}\\[4pt]
  {\large A plain-English guide to the n8n multi-agent workflow}\\[8pt]
  {\small\color{ink}\url{https://github.com/Subharjun/n8n-dependency-eol-firewall}}
\end{center}
\vspace{6pt}
\hrule
\vspace{14pt}

% ======================= 1. THE PROBLEM =======================
\section{The problem, in one minute}

Every software project depends on \textbf{other people's code} --- small building blocks called
\emph{dependencies} (for example \texttt{express} or \texttt{chalk}). These blocks keep changing:

\begin{itemize}[leftmargin=1.4em]
  \item Some get \textbf{deprecated} --- the author says ``stop using this.''
  \item Some reach \textbf{end-of-life (EOL)} --- no more security fixes.
  \item Some release a new \textbf{major version} with \textbf{breaking changes}.
\end{itemize}

If you don't notice in time, your project quietly becomes insecure or breaks on the next update.
Existing tools either spam you with blind ``bump the version'' pull requests, or list security holes
without telling you \emph{what to actually do}. \textbf{Nobody reads your real code and gives you a
judged, ready-to-act plan.} That is the gap this workflow fills.

% ======================= 2. THE IDEA =======================
\section{The idea, in one paragraph}

We built an automated ``firewall'' that runs on a schedule (say, every Monday). It looks at one GitHub
repository, checks each dependency against the official registry, and asks \textbf{three specialised AI
assistants (``agents'')} to cooperate: one decides \emph{whether} a dependency is genuinely risky,
one figures out \emph{how much of your code it affects}, and one \emph{writes the fix} as a real
GitHub Pull Request. On a healthy project it does nothing --- no noise.

\begin{center}
\colorbox{soft}{\parbox{0.92\linewidth}{\centering\small
\textbf{Guiding principle:} \emph{The AI judges; plain code edits.}\\
The agents only \emph{reason}. The actual version change in your file is done by exact,
deterministic code --- so the edit is always correct.}}
\end{center}

% ======================= 3. VOCABULARY =======================
\section{A tiny vocabulary (so nothing is confusing)}

\begin{description}[leftmargin=0em, style=nextline]
  \item[n8n] A visual ``automation'' tool. You connect boxes (called \textbf{nodes}) with arrows, and
    data flows left-to-right through them. Think of it as a flowchart that actually runs.
  \item[Node] One box = one action (fetch a web page, run a bit of code, call an AI model, \dots).
  \item[LLM / model] A large language model (here we use \textbf{Groq}'s \texttt{gpt-oss-120b}). It reads
    text and writes text.
  \item[Agent] An LLM node given a \textbf{specific job} and forced to answer in a strict, structured
    format (JSON) so the next node can use its answer reliably.
  \item[Pull Request (PR)] The standard way to propose a code change on GitHub, so a human can review
    and merge it.
\end{description}

% ======================= 4. THE BIG PICTURE =======================
\section{The big picture}

The workflow is a single chain of nodes. Data starts at the top and flows down. Purple boxes are the
three AI agents; blue boxes are ordinary automation steps.

\begin{center}
\begin{tikzpicture}[node distance=6mm]
  \node[step] (trig) {\textbf{1. Weekly Schedule} --- starts the workflow automatically.};
  \node[step, below=of trig] (cfg) {\textbf{2. Config: owner/repo} --- the \emph{one} place you type which
     GitHub repo to watch.};
  \node[step, below=of cfg] (read) {\textbf{3--6. Gather facts} --- read \texttt{package.json}, list every
     dependency, look each one up in the npm registry, compute ``how far behind / how stale.''};
  \node[ai, below=of read] (a1) {\textbf{7. Agent 1 -- Watcher} --- \emph{Is this genuinely risky?}
     Answers yes/no + a severity, conservatively.};
  \node[step, below=of a1] (filt) {\textbf{8. Filter} --- keep only the risky ones. If none, the workflow
     stops here (healthy repo).};
  \node[ai, below=of filt] (a2) {\textbf{9--11. Agent 2 -- Blast-Radius} --- searches \emph{your} code and
     reports which parts actually use this dependency and what breaks.};
  \node[ai, below=of a2] (a3) {\textbf{12. Agent 3 -- Migration Drafter} --- writes the upgrade plan:
     summary, risk, exact command, migration steps, and a diff.};
  \node[step, below=of a3] (pr) {\textbf{13--19. Open the Pull Request} --- detect the default branch,
     make a new branch, bump the version, commit, and open the PR.};

  \draw[arrow] (trig)--(cfg);
  \draw[arrow] (cfg)--(read);
  \draw[arrow] (read)--(a1);
  \draw[arrow] (a1)--(filt);
  \draw[arrow] (filt)--(a2);
  \draw[arrow] (a2)--(a3);
  \draw[arrow] (a3)--(pr);
\end{tikzpicture}
\end{center}

% ======================= 5. STEP BY STEP =======================
\section{Step-by-step: what every node does}

\subsection{Setup steps (gathering facts)}

\begin{enumerate}[leftmargin=1.5em]
  \item \textbf{Weekly Schedule.} A timer node. Every Monday at 06:00 it triggers a run. (You can also
    press ``Execute'' by hand any time.)
  \item \textbf{Config: owner/repo.} A tiny node holding two values, \texttt{owner} and \texttt{repo}.
    Every GitHub step reads the target repository from here, so you configure the workflow in exactly
    one place.
  \item \textbf{Read package.json.} Downloads the project's \texttt{package.json} file (the list of
    dependencies) from GitHub.
  \item \textbf{Parse Manifest.} A small code node that turns that one file into \emph{many} items ---
    one per dependency --- each carrying its name and installed version.
  \item \textbf{npm Registry Lookup.} For each dependency, it asks the public npm registry: what is the
    latest version? Is it deprecated? When was it last published?
  \item \textbf{Assess Registry Facts.} Turns the registry answer into hard numbers: how many major
    versions behind you are, how many days since the last release, any deprecation message. These
    ``facts'' are what the first agent will judge.
\end{enumerate}

\subsection{The decision (Agent 1)}

\begin{enumerate}[leftmargin=1.5em, start=7]
  \item \textbf{Agent 1 -- Watcher / Triage.} Reads the facts for one dependency and returns strict JSON
    like \texttt{\{at\_risk, risk\_type, severity, summary\}}. It is told to be \emph{conservative}: a
    slightly-old-but-fine package is \emph{not} flagged. This prevents spam.
  \item \textbf{Filter (``Only At-Risk'').} Lets through only the dependencies the Watcher flagged.
    Everything else stops. \emph{This is why a healthy repo produces zero pull requests.}
\end{enumerate}

\subsection{The investigation (Agent 2)}

\begin{enumerate}[leftmargin=1.5em, start=9]
  \item \textbf{Search Repo For Usage.} Uses GitHub's code search to find where in \emph{your} project
    the dependency is actually used.
  \item \textbf{Summarize Usage.} Compresses those search hits into a short briefing (file paths +
    snippets).
  \item \textbf{Agent 2 -- Blast-Radius.} Reads that briefing and reports which functions/symbols you
    use and how much an upgrade will disturb --- based only on evidence it was given, so it does not
    make things up.
\end{enumerate}

\subsection{The fix (Agent 3 + plain code)}

\begin{enumerate}[leftmargin=1.5em, start=12]
  \item \textbf{Agent 3 -- Migration Drafter.} Writes the human-facing content: a summary, a risk
    assessment, the exact upgrade command, a migration checklist, and a diff.
  \item \textbf{Prep PR Vars.} Prepares a unique branch name and the PR title/body.
  \item \textbf{Get Default Branch.} Asks GitHub whether the repo's main branch is called \texttt{main},
    \texttt{master}, or something else --- so it works on \emph{any} repository.
  \item \textbf{Get Base SHA \& Create Branch.} Makes a fresh branch to hold the change.
  \item \textbf{Bump package.json (deterministic).} Plain code changes the version number precisely
    (keeping symbols like \texttt{\^{}} or \texttt{\~{}}). \textbf{No AI touches the actual edit.}
  \item \textbf{Commit \& Open Pull Request.} Commits the bumped file and opens a Pull Request, with the
    agent's plan as the description, ready for a human to review and merge.
\end{enumerate}

% ======================= 6. WHY THREE AGENTS =======================
\section{Why three agents instead of one prompt}

\begin{center}
\begin{tabular}{@{}p{3.1cm}p{4.2cm}p{6.4cm}@{}}
\toprule
\textbf{Agent} & \textbf{Question it answers} & \textbf{Why it is separate} \\
\midrule
Watcher & ``Is this worth acting on?'' & A cheap, conservative gate so the expensive steps only run on
 real problems. \\
Blast-Radius & ``How much of \emph{my} code breaks?'' & Needs your repository's code as evidence; a
 different, focused task. \\
Drafter & ``Write the fix.'' & Turning a verdict into a reviewable PR is its own writing job. \\
\bottomrule
\end{tabular}
\end{center}

\noindent Splitting the work keeps each step simple, cheap, and reliable --- and lets the conservative
gate stop noise before any writing happens.

% ======================= 7. SETUP =======================
\section{How to run it yourself (5 minutes)}

\begin{enumerate}[leftmargin=1.5em]
  \item \textbf{Import} one of the two \texttt{.json} files into n8n
    (\emph{Workflows $\rightarrow$ Import from File}).
  \item \textbf{Add two credentials} and select them on the nodes: a \textbf{Groq} API key (free tier at
    \href{https://console.groq.com}{console.groq.com}) and a \textbf{GitHub} token with \texttt{repo}
    scope.
  \item \textbf{Open the ``Config: owner/repo'' node} and type your GitHub username and repository name.
  \item \textbf{Press ``Execute Workflow''} to test, then toggle it \textbf{Active} for the weekly run.
\end{enumerate}

\noindent The two versions:
\begin{itemize}[leftmargin=1.4em]
  \item \texttt{dependency-eol-firewall-autopr.json} --- opens a \textbf{Pull Request}.
  \item \texttt{dependency-eol-firewall.json} --- files a \textbf{GitHub Issue} instead (no code change).
\end{itemize}

% ======================= 8. OVERLEAF SYNC =======================
\section{Editing this document in Overleaf}

This file lives in the GitHub repo at \texttt{docs/explainer.tex}. To edit it in Overleaf and keep it
synced:

\begin{enumerate}[leftmargin=1.5em]
  \item In Overleaf: \emph{New Project $\rightarrow$ Import from GitHub} and pick
    \texttt{n8n-dependency-eol-firewall}. \emph{(Requires linking your GitHub account to Overleaf.)}
  \item Edit here; use Overleaf's \emph{Menu $\rightarrow$ GitHub $\rightarrow$ Push} to send changes back,
    and \emph{Pull} to fetch updates.
  \item If you'd rather not link accounts: download this \texttt{.tex} and drag it into a blank Overleaf
    project. It compiles with the default \textbf{pdfLaTeX} setting --- no extra configuration.
\end{enumerate}

\vspace{10pt}
\hrule
\vspace{6pt}
\begin{center}
\small\color{ink} Built with n8n + Groq. MIT licensed. Judge with AI, edit with code.
\end{center}

\end{document}
